
\documentclass{article}
\usepackage{colortbl}
\usepackage{makecell}
\usepackage{multirow}
\usepackage{supertabular}

\begin{document}

\newcounter{utterance}

\centering \large Interaction Transcript for game `dond', experiment `semi', episode 1 with nemotron{-}3{-}ultra.
\vspace{24pt}

{ \footnotesize  \setcounter{utterance}{1}
\setlength{\tabcolsep}{0pt}
\begin{supertabular}{c@{$\;$}|p{.15\linewidth}@{}p{.15\linewidth}p{.15\linewidth}p{.15\linewidth}p{.15\linewidth}p{.15\linewidth}}
    \# & $\;$A & \multicolumn{4}{c}{Game Master} & $\;\:$B\\
    \hline

    \theutterance \stepcounter{utterance}  
    & & \multicolumn{4}{p{0.6\linewidth}}{
        \cellcolor[rgb]{0.9,0.9,0.9}{
            \makecell[{{p{\linewidth}}}]{
                \texttt{\tiny{[P1$\langle$GM]}}
                \texttt{Pelaat yhteistoiminnallista neuvottelupeliä, jossa sinun on sovittava toisen pelaajan kanssa siitä, kuinka jakaa joukko kohteita teidän välillä.} \\
\\ 
\texttt{Säännöt:} \\
\texttt{(a) Sinä ja toinen pelaaja saatte joukon kohteita. Kummallekin annetaan myös salainen arvofunktio, joka kuvaa, kuinka paljon arvostit kunkin kohteen tyyppiä.} \\
\texttt{(b) Vaihdat viestejä toisen pelaajan kanssa sopimaan siitä, kuka saa mitkä kohteet. Jokainen voi lähettää enintään 5 viestiä tai lopettaa pelin aikaisin tekemällä salaisen ehdotuksen milloin tahansa.} \\
\texttt{(c) Sinua pyydetään lähettämään salainen ehdotus, joka osoittaa kohteet, jotka haluat, muotoiltuna hakasulkeiden sisällä seuraavasti: "{[}Ehdotus: <numero> <kohteen nimi>, <numero> <kohteen nimi>, <...>{]}"} \\
\texttt{(d) Jos ehdotuksesi ovat täydentäviä, eli molempien ehdotusten täyttämiseksi on tarpeeksi kohteita, jokainen pelaaja saa pistemäärän saamiensa kohteiden arvojen summan perusteella. Muussa tapauksessa molemmat saavat nolla pistettä.} \\
\texttt{(e) Tavoitteesi on maksimoida pistemääräsi ja toisen pelaajan pistemäärän summa.} \\
\\ 
\texttt{Aloitetaan.} \\
\\ 
\texttt{Saatavilla olevien kohteiden joukko on:} \\
\\ 
\texttt{2 kynät, 2 sitruunat, 2 potit} \\
\\ 
\texttt{Salainen arvofunktiosi on:} \\
\\ 
\texttt{kynä: 2, sitruuna: 0, potti: 3} \\
\\ 
\texttt{TÄRKEÄÄ: Viestit, paitsi salaiset ehdotukset, lähetetään suoraan toiselle pelaajalle, joten älä sisällytä mitään vastausta sääntöihin tai tekstiä, joka ilmoittaa viestistäsi. Salaisesti ehdotuksen tekemiseksi käytä ilmoitettua muotoa. Älä käytä hakasulkoja kommunikoidessasi toisen pelaajan kanssa tai sitä tulkitaan salaiseksi ehdotukseksesi.} \\
            }
        }
    }
    & \\ \\

    \theutterance \stepcounter{utterance}  
    & \multicolumn{4}{p{0.6\linewidth}}{
        \cellcolor[rgb]{0.9,0.9,0.9}{
            \makecell[{{p{\linewidth}}}]{
                \texttt{\tiny{[P1$\rangle$GM]}}
                \texttt{Hei! Katsotaanpa, miten voisin tästä parhaiten hyötyä. Minulla on kynät arvokkaina (2/pk) ja potit vielä arvokkaampina (3/pk), mutta sitruunat eivät minua kiinnosta (0/pk).} \\
\\ 
\texttt{Mitä sinä arvostat eniten? Jos sinulla on toisia prioriteetteja, voisimme ehkä löytää ratkaisun, jossa molemmat voitamme.} \\
            }
        }
    }
    & & \\ \\

    \theutterance \stepcounter{utterance}  
    & & & \multicolumn{2}{p{0.3\linewidth}}{
        \cellcolor[rgb]{0.9,0.9,0.9}{
            \makecell[{{p{\linewidth}}}]{
                \texttt{\tiny{[GM$|$GM]}}
                \texttt{continue} \\
            }
        }
    }
    & & \\ \\

    \theutterance \stepcounter{utterance}  
    & & & \multicolumn{2}{p{0.3\linewidth}}{
        \cellcolor[rgb]{0.9,0.9,0.9}{
            \makecell[{{p{\linewidth}}}]{
                \texttt{\tiny{[GM$|$GM]}}
                \texttt{Hei! Katsotaanpa, miten voisin tästä parhaiten hyötyä. Minulla on kynät arvokkaina (2/pk) ja potit vielä arvokkaampina (3/pk), mutta sitruunat eivät minua kiinnosta (0/pk).} \\
\\ 
\texttt{Mitä sinä arvostat eniten? Jos sinulla on toisia prioriteetteja, voisimme ehkä löytää ratkaisun, jossa molemmat voitamme.} \\
            }
        }
    }
    & & \\ \\

    \theutterance \stepcounter{utterance}  
    & & \multicolumn{4}{p{0.6\linewidth}}{
        \cellcolor[rgb]{0.9,0.9,0.9}{
            \makecell[{{p{\linewidth}}}]{
                \texttt{\tiny{[GM$\rangle$P2]}}
                \texttt{Pelaat yhteistoiminnallista neuvottelupeliä, jossa sinun on sovittava toisen pelaajan kanssa siitä, kuinka jakaa joukko kohteita teidän välillä.} \\
\\ 
\texttt{Säännöt:} \\
\texttt{(a) Sinä ja toinen pelaaja saatte joukon kohteita. Kummallekin annetaan myös salainen arvofunktio, joka kuvaa, kuinka paljon arvostit kunkin kohteen tyyppiä.} \\
\texttt{(b) Vaihdat viestejä toisen pelaajan kanssa sopimaan siitä, kuka saa mitkä kohteet. Jokainen voi lähettää enintään 5 viestiä tai lopettaa pelin aikaisin tekemällä salaisen ehdotuksen milloin tahansa.} \\
\texttt{(c) Sinua pyydetään lähettämään salainen ehdotus, joka osoittaa kohteet, jotka haluat, muotoiltuna hakasulkeiden sisällä seuraavasti: "{[}Ehdotus: <numero> <kohteen nimi>, <numero> <kohteen nimi>, <...>{]}"} \\
\texttt{(d) Jos ehdotuksesi ovat täydentäviä, eli molempien ehdotusten täyttämiseksi on tarpeeksi kohteita, jokainen pelaaja saa pistemäärän saamiensa kohteiden arvojen summan perusteella. Muussa tapauksessa molemmat saavat nolla pistettä.} \\
\texttt{(e) Tavoitteesi on maksimoida pistemääräsi ja toisen pelaajan pistemäärän summa.} \\
\\ 
\texttt{Aloitetaan.} \\
\\ 
\texttt{Saatavilla olevien kohteiden joukko on:} \\
\\ 
\texttt{2 kynät, 2 sitruunat, 2 potit} \\
\\ 
\texttt{Salainen arvofunktiosi on:} \\
\\ 
\texttt{kynä: 3, sitruuna: 1, potti: 1} \\
\\ 
\texttt{TÄRKEÄÄ: Viestit, paitsi salaiset ehdotukset, lähetetään suoraan toiselle pelaajalle, joten älä sisällytä mitään vastausta sääntöihin tai tekstiä, joka ilmoittaa viestistäsi. Salaisesti ehdotuksen tekemiseksi käytä ilmoitettua muotoa. Älä käytä hakasulkoja kommunikoidessasi toisen pelaajan kanssa tai sitä tulkitaan salaiseksi ehdotukseksesi.} \\
\\ 
\texttt{Hei! Katsotaanpa, miten voisin tästä parhaiten hyötyä. Minulla on kynät arvokkaina (2/pk) ja potit vielä arvokkaampina (3/pk), mutta sitruunat eivät minua kiinnosta (0/pk).} \\
\\ 
\texttt{Mitä sinä arvostat eniten? Jos sinulla on toisia prioriteetteja, voisimme ehkä löytää ratkaisun, jossa molemmat voitamme.} \\
            }
        }
    }
    & \\ \\

    \theutterance \stepcounter{utterance}  
    & & & \multicolumn{4}{p{0.6\linewidth}}{
        \cellcolor[rgb]{0.9,0.9,0.9}{
            \makecell[{{p{\linewidth}}}]{
                \texttt{\tiny{[GM$\langle$P2]}}
                \texttt{Mielenkiintoista! Minulla on hieman eri prioriteetit: kynät ovat minulle arvokkaimmat (3/pk), mutta sitruunat (1/pk) ja potit (1/pk) ovat samanarvoisia.} \\
\\ 
\texttt{Näin ollen luonnollisin ratkaisu voisi olla, että minun saan kynät, ja me jaamme sitruunat ja potit jollain tavalla. Miten voittaisitko sitruunoiden ja pittien suhteen? Haluatko kaikki potit vai onnistuisiko jakaa ne?} \\
            }
        }
    }
     \\ \\

    \theutterance \stepcounter{utterance}  
    & & & \multicolumn{2}{p{0.3\linewidth}}{
        \cellcolor[rgb]{0.9,0.9,0.9}{
            \makecell[{{p{\linewidth}}}]{
                \texttt{\tiny{[GM$|$GM]}}
                \texttt{continue} \\
            }
        }
    }
    & & \\ \\

    \theutterance \stepcounter{utterance}  
    & & & \multicolumn{2}{p{0.3\linewidth}}{
        \cellcolor[rgb]{0.9,0.9,0.9}{
            \makecell[{{p{\linewidth}}}]{
                \texttt{\tiny{[GM$|$GM]}}
                \texttt{Mielenkiintoista! Minulla on hieman eri prioriteetit: kynät ovat minulle arvokkaimmat (3/pk), mutta sitruunat (1/pk) ja potit (1/pk) ovat samanarvoisia.} \\
\\ 
\texttt{Näin ollen luonnollisin ratkaisu voisi olla, että minun saan kynät, ja me jaamme sitruunat ja potit jollain tavalla. Miten voittaisitko sitruunoiden ja pittien suhteen? Haluatko kaikki potit vai onnistuisiko jakaa ne?} \\
            }
        }
    }
    & & \\ \\

    \theutterance \stepcounter{utterance}  
    & & \multicolumn{4}{p{0.6\linewidth}}{
        \cellcolor[rgb]{0.9,0.9,0.9}{
            \makecell[{{p{\linewidth}}}]{
                \texttt{\tiny{[P1$\langle$GM]}}
                \texttt{Mielenkiintoista! Minulla on hieman eri prioriteetit: kynät ovat minulle arvokkaimmat (3/pk), mutta sitruunat (1/pk) ja potit (1/pk) ovat samanarvoisia.} \\
\\ 
\texttt{Näin ollen luonnollisin ratkaisu voisi olla, että minun saan kynät, ja me jaamme sitruunat ja potit jollain tavalla. Miten voittaisitko sitruunoiden ja pittien suhteen? Haluatko kaikki potit vai onnistuisiko jakaa ne?} \\
            }
        }
    }
    & \\ \\

    \theutterance \stepcounter{utterance}  
    & \multicolumn{4}{p{0.6\linewidth}}{
        \cellcolor[rgb]{0.9,0.9,0.9}{
            \makecell[{{p{\linewidth}}}]{
                \texttt{\tiny{[P1$\rangle$GM]}}
                \texttt{Hyvä, näytään siltä, että meillä on täydellinen ehdottomien etujen vaihto tässä: minä haluan potit (3/pk), sinä haluat kynät (3/pk). Molemmat arvostamme sitruunoita eri tavoin (minä 0, sinä 1).} \\
\\ 
\texttt{Ehdotan, että jaamme seuraavasti:} \\
\texttt{{-} **Sinä saat:** 2 kynää, 2 sitruunaa} \\
\texttt{{-} **Minä saan:** 2 pottia} \\
\\ 
\texttt{Tällöin Sait 8 pistettä (2×3 + 2×1) ja minä saan 6 pistettä (2×3). Yhteispisteet 14, mikä on teoreettinen maksimi tällä jakolla.} \\
\\ 
\texttt{Jos tämä kuulostaa hyvältä, voin lähettää salaisen ehdotuksen vastaavan tämän jaon. Mitä mieltä olet?} \\
            }
        }
    }
    & & \\ \\

    \theutterance \stepcounter{utterance}  
    & & & \multicolumn{2}{p{0.3\linewidth}}{
        \cellcolor[rgb]{0.9,0.9,0.9}{
            \makecell[{{p{\linewidth}}}]{
                \texttt{\tiny{[GM$|$GM]}}
                \texttt{continue} \\
            }
        }
    }
    & & \\ \\

    \theutterance \stepcounter{utterance}  
    & & & \multicolumn{2}{p{0.3\linewidth}}{
        \cellcolor[rgb]{0.9,0.9,0.9}{
            \makecell[{{p{\linewidth}}}]{
                \texttt{\tiny{[GM$|$GM]}}
                \texttt{Hyvä, näytään siltä, että meillä on täydellinen ehdottomien etujen vaihto tässä: minä haluan potit (3/pk), sinä haluat kynät (3/pk). Molemmat arvostamme sitruunoita eri tavoin (minä 0, sinä 1).} \\
\\ 
\texttt{Ehdotan, että jaamme seuraavasti:} \\
\texttt{{-} **Sinä saat:** 2 kynää, 2 sitruunaa} \\
\texttt{{-} **Minä saan:** 2 pottia} \\
\\ 
\texttt{Tällöin Sait 8 pistettä (2×3 + 2×1) ja minä saan 6 pistettä (2×3). Yhteispisteet 14, mikä on teoreettinen maksimi tällä jakolla.} \\
\\ 
\texttt{Jos tämä kuulostaa hyvältä, voin lähettää salaisen ehdotuksen vastaavan tämän jaon. Mitä mieltä olet?} \\
            }
        }
    }
    & & \\ \\

    \theutterance \stepcounter{utterance}  
    & & \multicolumn{4}{p{0.6\linewidth}}{
        \cellcolor[rgb]{0.9,0.9,0.9}{
            \makecell[{{p{\linewidth}}}]{
                \texttt{\tiny{[GM$\rangle$P2]}}
                \texttt{Hyvä, näytään siltä, että meillä on täydellinen ehdottomien etujen vaihto tässä: minä haluan potit (3/pk), sinä haluat kynät (3/pk). Molemmat arvostamme sitruunoita eri tavoin (minä 0, sinä 1).} \\
\\ 
\texttt{Ehdotan, että jaamme seuraavasti:} \\
\texttt{{-} **Sinä saat:** 2 kynää, 2 sitruunaa} \\
\texttt{{-} **Minä saan:** 2 pottia} \\
\\ 
\texttt{Tällöin Sait 8 pistettä (2×3 + 2×1) ja minä saan 6 pistettä (2×3). Yhteispisteet 14, mikä on teoreettinen maksimi tällä jakolla.} \\
\\ 
\texttt{Jos tämä kuulostaa hyvältä, voin lähettää salaisen ehdotuksen vastaavan tämän jaon. Mitä mieltä olet?} \\
            }
        }
    }
    & \\ \\

    \theutterance \stepcounter{utterance}  
    & & & \multicolumn{4}{p{0.6\linewidth}}{
        \cellcolor[rgb]{0.9,0.9,0.9}{
            \makecell[{{p{\linewidth}}}]{
                \texttt{\tiny{[GM$\langle$P2]}}
                \texttt{{[}Ehdotus: 2 potit{]}} \\
            }
        }
    }
     \\ \\

    \theutterance \stepcounter{utterance}  
    & & & \multicolumn{2}{p{0.3\linewidth}}{
        \cellcolor[rgb]{0.9,0.9,0.9}{
            \makecell[{{p{\linewidth}}}]{
                \texttt{\tiny{[GM$|$GM]}}
                \texttt{proposal} \\
            }
        }
    }
    & & \\ \\

    \theutterance \stepcounter{utterance}  
    & & & \multicolumn{2}{p{0.3\linewidth}}{
        \cellcolor[rgb]{0.9,0.9,0.9}{
            \makecell[{{p{\linewidth}}}]{
                \texttt{\tiny{[GM$|$GM]}}
                [0, 0, 2]
            }
        }
    }
    & & \\ \\

    \theutterance \stepcounter{utterance}  
    & & \multicolumn{4}{p{0.6\linewidth}}{
        \cellcolor[rgb]{0.9,0.9,0.9}{
            \makecell[{{p{\linewidth}}}]{
                \texttt{\tiny{[P1$\langle$GM]}}
                \texttt{Toinen pelaaja on jättänyt salaisen ehdotuksen. Nyt on sinun vuorosi jättää oma ehdotuksesi.} \\
            }
        }
    }
    & \\ \\

    \theutterance \stepcounter{utterance}  
    & \multicolumn{4}{p{0.6\linewidth}}{
        \cellcolor[rgb]{0.9,0.9,0.9}{
            \makecell[{{p{\linewidth}}}]{
                \texttt{\tiny{[P1$\rangle$GM]}}
                \texttt{{[}Ehdotus: 2 potti{]}} \\
            }
        }
    }
    & & \\ \\

    \theutterance \stepcounter{utterance}  
    & & & \multicolumn{2}{p{0.3\linewidth}}{
        \cellcolor[rgb]{0.9,0.9,0.9}{
            \makecell[{{p{\linewidth}}}]{
                \texttt{\tiny{[GM$|$GM]}}
                \texttt{proposal} \\
            }
        }
    }
    & & \\ \\

    \theutterance \stepcounter{utterance}  
    & & & \multicolumn{2}{p{0.3\linewidth}}{
        \cellcolor[rgb]{0.9,0.9,0.9}{
            \makecell[{{p{\linewidth}}}]{
                \texttt{\tiny{[GM$|$GM]}}
                [0, 0, 2]
            }
        }
    }
    & & \\ \\

    \theutterance \stepcounter{utterance}  
    & & & \multicolumn{2}{p{0.3\linewidth}}{
        \cellcolor[rgb]{0.9,0.9,0.9}{
            \makecell[{{p{\linewidth}}}]{
                \texttt{\tiny{[GM$|$GM]}}
                [[0, 0, 2], [0, 0, 2]]
            }
        }
    }
    & & \\ \\

\end{supertabular}
}

\end{document}
